Illustrations can be included using the command \lstinline+\includegraphics+
from the package~\ctanref{graphicx}. The image must be available as a graphic
file for this. Place illustrations in the subdirectory \texttt{figures}.
(see section~\ref{sec:tpl:dir}). Depending on the application, use the
following image formats.

\begin{description}
\item[PDF] Technical drawings, graphics, diagrams, etc. -- In general
  so-called vector graphics are meant here. These consist of points, lines
  and areas and can be enlarged or reduced without loss of quality.
  They also take up less storage space.

  Please note, however, that texts within vector graphics are then also
  scaled as well and this may lead to an inconsistent text image.

\item[PNG] Screenshots, computer-generated images (3D visualisations,
  simulation results, \dots), land or weather maps, heatmaps -- These are so-called
  raster graphics. Such images consist of pixels, which are coloured
  differently. Since the pixels are not perceived individually, the result is
  an image impression. The enlargement or reduction of raster images, however,
  leads to loss of quality (blurring, loss of detail). Raster graphics
  require more storage space than vector graphics.

  PNG files are a lossless image format, i.e. the image information is not lost
  during storage. It is therefore for images with fine details (e.g. lines) and
  high-contrast images (black and white drawings).

\item[JPEG] Photographs, scanned images -- These are also raster graphics.
  However, a lossy storage method is used. Colour values are changed by this.
  At natural way'', such as photographs or scanned documents, these changes are
  not noticeable. Unnatural images with strong contrasts, fine lines or
  strongly saturated colours, however, lead to display errors
  (so-called artefacts), which is why the JPEG format should not be used in
  this case. Due to the storage technology, JPEG images generally require less
  memory than PNG images.
\end{description}

Images should always be placed within a \texttt{figure} environment. This ensures that
the illustrations are numbered and that they are inclusion in the list of
figures. An example can be found in Listing~\ref{lst:latex:figure}.

\begin{lstlisting}[language={[LaTeX]{TeX}},style=block,escapechar=\!,float,caption={An illustration},label={lst:latex:figure}]
\begin{figure}
\centering
\includegraphics[width=0.7\linewidth]{figures/!\lara{filename}!}
\caption{!\lara{caption}!}
\label{!\lara{reference mark}!}
\end{figure}
\end{lstlisting}

No file extension needs to be -- and should be -- specified in the file name. The
reference mark can be omitted if no cross-reference to the image is planned.

The \texttt{figure} environment is a sliding environment. The figure can be moved
automatically to another place in the document if space availability requires
it. For more information on this see section~\ref{sec:final:float}.
